\section*{Conclusion}

This study represents a systematic empirical validation of the ConvMixer architecture in the audio domain, providing evidence for its applicability to lightweight speech recognition systems in edge-AI environments. By adapting a vision-inspired patch-based convolutional framework to Mel-spectrogram representations, we show that ConvMixer can achieve high accuracy with low computational cost, effectively bridging vision-derived patch design and efficient speech understanding.

The proposed system, featuring an optimized preprocessing pipeline (denoising, WebRTC VAD \cite{WebRTCVAD}, normalization) and a compact ConvMixer-256/8 configuration (0.7M parameters), achieved representative single-run performance of 97.42\% test accuracy and mean F1-score of 0.97 across 13 Vietnamese speech command classes. The revised manuscript further introduces a controlled benchmark protocol across ConvMixer-256/8, ResNet-18, MobileNetV2, EfficientNet, and AST (tiny configuration) with five seeds (42--46), enabling fairer and more statistically reliable comparisons.

Overall, this research supports ConvMixer as a resource-efficient and competitive option for real-time speech command recognition under controlled settings. However, the current study remains limited by dataset size, language scope, and benchmark breadth. Future work will extend controlled comparisons to additional transformer variants, explore large-scale pretraining and hybrid ConvMixer–attention designs, and optimize inference for embedded and IoT deployment, while broadening evaluation to non-speech audio tasks.
